\section{Resumen}

Se evaluó el efecto de campos magnéticos estáticos sobre la germinación y el crecimiento inicial de semillas de \textit{Zea mays}. Se realizaron tres controles negativos de 72\,h cada uno (\(n = 81\) bandejas en total) para caracterizar la variabilidad espacial del sistema experimental, seguidos de experimentos con imanes cerámicos de ferrita en configuración polo norte, exponiendo a las semillas durante las primeras 72\,h de germinación. La variable principal de análisis fue el cambio porcentual en peso húmedo.

En los controles negativos se detectó un efecto espacial asociado al piso de la bandeja (A \(>\) B \(>\) C, \(p = 0.0228\)), pero no a las posiciones laterales ni de profundidad. El análisis a escala de celda individual mostró que las diferencias estuvieron impulsadas principalmente por la celda C9, que presentó menor crecimiento. Se verificó que no hubo sesgo de selección en el armado de las bandejas (\(p = 0.4926\)).

El estudio buscó identificar condiciones asociadas a mejoras en variables de crecimiento temprano, priorizando un diseño experimental sistemático, controlado y cuantitativo. Se estimaron tamaño de efecto, significancia estadística y reproducibilidad de los resultados bajo condiciones controladas de germinación.
